\documentclass[12pt,a4paper]{article}

\usepackage[utf8]{inputenc}
\usepackage[T1]{fontenc}
\usepackage{lmodern}
\usepackage{amsmath,amsthm,amssymb,amsfonts}
\usepackage{mathtools}
\usepackage{hyperref}
\usepackage{booktabs}
\usepackage[margin=1in]{geometry}

\newtheorem{theorem}{Theorem}[section]
\newtheorem{lemma}[theorem]{Lemma}
\newtheorem{proposition}[theorem]{Proposition}
\newtheorem{corollary}[theorem]{Corollary}
\newtheorem{definition}[theorem]{Definition}
\newtheorem{remark}[theorem]{Remark}

\newcommand{\N}{\mathbb{N}}
\newcommand{\Z}{\mathbb{Z}}
\newcommand{\R}{\mathbb{R}}
\newcommand{\C}{\mathbb{C}}
\newcommand{\Syr}{\mathrm{Syr}}
\newcommand{\vtwo}{v_2}
\newcommand{\cert}{\mathrm{cert}}
\newcommand{\tril}{\mathrm{tril}}
\newcommand{\Sodd}{S^{\mathrm{odd}}}

\title{An Elementary Uniform Spectral Gap\\for the Syracuse Transfer Operator}

\author{Macdara Ó Murchú\thanks{With Claude (Anthropic): Opus~4.8 research sessions and a
Fable~5 adversarial review, refereeing, and consolidation pass. All proofs are elementary and
every numerical claim is reproducible from the repository.}\\
\small \url{https://github.com/m4cd4r4/collatz-conjecture-spectral-gap}}

\date{July 2026}

\begin{document}

\maketitle

\begin{abstract}
Let $T_k$ be the transfer operator of the Syracuse map $\Syr(n) = (3n+1)/2^{\vtwo(3n+1)}$ on the
$N = 2^{k-1}$ odd residues modulo $2^k$, averaging over the $2^k$ lifts of a residue. We prove a
uniform spectral gap: for every $k \geq 3$,
\[
  |\lambda_2(T_k)| \;\leq\; \cert(k) \;\leq\; 2^{-3/2} + 2^{-1} \;=\; 0.8535\ldots \;<\; 1,
\]
where $\cert(k)$ is an explicit weighted row-sum certificate over a 2-adic frequency grading.
The proof is entirely elementary: finite cyclic 2-group arithmetic (three uniform-fiber lemmas
sharing one engine, ``3 is a unit modulo $2^k$''), a finite geometric series, an integer collision
count, and a single application of the Cauchy--Schwarz inequality. A sharper per-level bound on
the defect covector improves the constant to $0.6827$; a further sharpening to $0.6553$ rests on
a constant verified only for $k \leq 26$ and is labelled as such throughout. The measured values
are $\cert(k) \approx 0.6345$ and $|\lambda_2| \approx 0.27$. The elementary core of the proof is
machine-checked in Lean~4 (Mathlib, sorry-free).

\textbf{What it measures.} The honest 2-adic kernel, conditioned on the full fibre rather than
on a finite lift window, is nilpotent off its Perron eigenvalue, so the mixing of the averaged
mod-$2^k$ dynamics is trivial. $T_k$ differs from it by a rank-one but \emph{not small} defect
($\|T_k - K_k\|_1 \approx 0.4$), and the theorem is the statement that this large rank-one
truncation defect cannot drive the spectrum to the unit circle at any scale --- a uniform
stability statement for the finite-window approximation.

\textbf{Scope.} This is a theorem about the Markov model, not about Collatz orbits. In
particular it does \emph{not} rule out Collatz cycles: the analogous $3x{-}1$ operator passes the
identical certificate even more strongly, yet $3x-1$ has non-trivial cycles. An earlier version
of this project claimed a cycle-elimination corollary; that claim was withdrawn (Section~\ref{sec:scope}).
\end{abstract}

\noindent\textbf{Keywords:} Syracuse map, transfer operator, spectral gap, 2-adic analysis.

%==============================================================================
\section{Introduction}
%==============================================================================

For odd $n$, the Syracuse map is $\Syr(n) = (3n+1)/2^{\vtwo(3n+1)}$, where $\vtwo$ is the 2-adic
valuation. Reducing modulo $2^k$ gives a natural Markov chain on the $N = 2^{k-1}$ odd residues:
from a residue $r$, move to $\Syr(r + m 2^k) \bmod 2^k$ with $m$ uniform on $[0, 2^k)$. Its
transition matrix $T_k$ (column-stochastic) is the finite-scale transfer operator of the map;
the mixing of this family is the standard probabilistic model for the equidistribution of
Syracuse trajectories modulo powers of two, going back to the heuristics surveyed in
\cite{Lagarias2010} and quantified in Tao's almost-all result \cite{Tao2019}.

Our main result is a uniform, explicit, and elementarily proved spectral gap for this family.

\begin{theorem}[Uniform spectral gap]\label{thm:main}
For every $k \geq 3$,
\[
  |\lambda_2(T_k)| \;\leq\; \cert(k) \;:=\; \max_{0 \le a \le k-2} \sum_{b=0}^{k-2} Q_k[a,b]\, 2^{a-b}
  \;\leq\; 2^{-3/2} + 2^{-1} \;<\; 1,
\]
where $\lambda_2$ is the second-largest eigenvalue in modulus and
$Q_k[a,b] = \|P_a U_k P_b\|_2$ is the matrix of block operator norms of the Koopman adjoint
$U_k = T_k^{\!\top}$ over the 2-adic frequency grading of Section~\ref{sec:setup}.
\end{theorem}

\subsection*{What the theorem says, informally}

A reader who wants the shape of the result before the notation may take the following as the
whole story.

The Collatz map is deterministic, and deterministic questions about it are hard. A standard
softening is to stop following one number and follow instead \emph{where numbers land on
average}, recording only the remainder modulo a power of two. That turns the map into a random
walk on a finite set of states --- the $2^{k-1}$ odd residues mod $2^k$ --- and $T_k$ is the
matrix of that walk. The averaging is over the $2^k$ integers sharing a given residue.

A random walk ``mixes'' when the memory of where it started decays. The rate is governed by the
second-largest eigenvalue $\lambda_2$: mixing happens exactly when $|\lambda_2| < 1$, and faster
the smaller it is. For a \emph{single} $k$ this is a finite computation and carries no
information. The content of Theorem~\ref{thm:main} is that one constant works for
\emph{every} $k$ at once, even though the matrix grows exponentially.

The proof has one idea, applied three times. Multiplication by $3$ is invertible modulo any power
of two --- because $3$ is odd --- so it permutes residues and spreads any set of them perfectly
evenly across the classes it can reach. Everything else is bookkeeping: sorting residues by how
divisible they are by two, summing a geometric series over those layers, counting collisions in
one exceptional row, and applying Cauchy--Schwarz once. There is no analysis and no computer
search inside the proof.

The one place the argument could have failed, and does not, is that top layer. The cascade of
blocks \emph{truncates} there rather than continuing, and it is exactly that truncation which
makes an $L^2$ estimate strong enough. Without it the bound would exceed $1$ and say nothing.

Finally, and this is not a caveat but the main thing to carry away: none of this bears on the
Collatz conjecture, and Section~\ref{sec:scope} shows by an explicit control that it cannot.
Averaging destroys the very structure a cycle lives in.

\subsection*{Why the theorem is of interest}

The interest of the theorem is threefold. First, uniformity: the bound is independent of $k$,
while the state space grows as $2^{k-1}$. Second, the constant is explicit and the true values
sit far below it ($|\lambda_2| \approx 0.27$ for all tested $k$), so the certificate has a large
margin. Third, and mainly, the proof is elementary from end to end. What looks like it should
require Gauss-sum estimates reduces to three instances of one fact (multiplication by a unit has
uniform fibers on a cyclic 2-group), one finite geometric series, one integer collision count
with a two-line injectivity argument, and one Cauchy--Schwarz application. No analytic number
theory, no perturbation theory, and no computer-assisted spectral certification enters the proof;
numerics appear only as consistency checks.

\label{sec:intro-history}%
We are explicit about history because this project previously overclaimed twice, and both
corrections shaped the final form. (i) A March 2026 draft asserted the gap via an eigenvalue
perturbation telescope with a ``bounded condition number''; an adversarial audit (Session 39)
refuted this: the eigenvector condition number grows like $10^k$, and the published inductive
inequality fails numerically at every odd $k \geq 7$. The present proof replaces eigenvalue
perturbation with norm inequalities that are condition-number-free. (ii) The certificate was
then presented as eliminating non-trivial Collatz cycles. That inference is false and was
withdrawn; Section~\ref{sec:scope} records the decisive control experiment. What survives is the
operator theorem itself.

\subsection*{Proof architecture}

\begin{enumerate}
\item \textbf{Spectral reduction} (Section~\ref{sec:reduction}): $|\lambda_2(T_k)| \le \rho(Q_k)
  \le \cert(k)$, by a one-sided Perron split, a level-energy majorisation applied to a
  compression, and Gelfand's formula. A subtlety repaired late in the project: $T_k$ is
  \emph{not} doubly stochastic and its stationary distribution is \emph{not} exactly uniform,
  so the argument must (and does) avoid both.
\item \textbf{Block structure} (Section~\ref{sec:blocks}): $U_k = U_{\mathrm{clean}} + D$ where
  $U_{\mathrm{clean}}$ is strictly upper-triangular in the level grading with exactly computable
  block norms $\|P_a U_{\mathrm{clean}} P_b\| = 2^{-(b-a)/2}$ (Lemma A), and $D$ is a rank-1
  defect supported at the single residue $r^\ast = -3^{-1} \bmod 2^k$.
\item \textbf{Defect mass} (Section~\ref{sec:lemmaB}): the defect covector $c$ satisfies
  $\|c\|_2 \le \sqrt{3}\,2^{-k/2}$, by an integer collision count (Lemma B).
\item \textbf{Assembly} (Section~\ref{sec:assembly}): each weighted row sum is bounded by the
  envelope $f(e) = \sum_{d=1}^{e-2} 2^{-3d/2} + 2^{-(e-1)/2}$ with $e = k - a$, whose maximum is
  $f(3) = 2^{-3/2} + 2^{-1}$. The only analytic input is Cauchy--Schwarz; crucially, the upper
  cascade \emph{truncates} at the top rows, which is what makes the $L^2$ bound sufficient.
\end{enumerate}

Section~\ref{sec:lemmaC} records the sharper per-level defect bound (Lemma C, via an elementary
frequency-doubling identity), which improves the constant to $0.6827$ but is not needed for
Theorem~\ref{thm:main}. Its sharp form gives $0.6553$ but consumes a constant verified only to
$k \leq 26$; we keep the two apart and never quote the sharper number as proved. Section~\ref{sec:lean} describes the Lean~4 formalisation of the
elementary core. Section~\ref{sec:scope} delimits scope.

%==============================================================================
\section{The operator and the frequency grading}\label{sec:setup}
%==============================================================================

\begin{definition}\label{def:T}
Fix $k \ge 3$ and let $N = 2^{k-1}$. The state space is the set of odd residues modulo $2^k$.
The chain moves from $r$ to $\Syr(r + m 2^k) \bmod 2^k$ with $m$ uniform on $\{0, \dots, 2^k-1\}$;
$T_k[s', s]$ is this transition probability, and $U = U_k := T_k^{\!\top}$ is the Koopman
(expectation) operator, $(Ug)(r) = \mathbb{E}_m\, g\!\left(\Syr(r + m 2^k)\right)$.
\end{definition}

$T_k$ is column-stochastic by construction, so $\mathbf{1}^{\!\top} T_k = \mathbf{1}^{\!\top}$
and $\rho(T_k) = 1$.

\begin{remark}[$T_k$ is not doubly stochastic]\label{rem:notds}
The row sums of $T_k$ are \emph{not} identically 1: numerically
$\|T_k \mathbf{1} - \mathbf{1}\|_2 \approx 0.764 \cdot 2^{-k/2}$, another face of the $r^\ast$
defect of Section~\ref{sec:blocks}. Consequently the stationary distribution is only
approximately uniform, $U$ does not preserve the mean-zero subspace, and identities such as
$T_r^p = T_k^p - N^{-1} J$ (used in an earlier write-up) are false. Nothing below uses
uniformity of the stationary distribution; the reduction in Section~\ref{sec:reduction} is
deliberately one-sided.
\end{remark}

\begin{definition}[Characters and levels]\label{def:levels}
Let $\omega = e^{2\pi i / 2^k}$ and, for $\xi \in [0, 2^{k-1})$, let
$\chi_\xi(r) = \omega^{\xi r}/\sqrt{N}$ on odd $r$. On odd arguments
$\chi_{\xi + 2^{k-1}} = -\chi_\xi$, so this range selects one representative per proportionality
class, and the $\chi_\xi$ form an orthonormal basis of $\C^N$ (orthogonality is the $m = k$ case
of Lemma~\ref{lem:S4} below). $\chi_0$ is the constant (Perron) vector. The \emph{level} of
$\xi \neq 0$ is $\vtwo(\xi) \in \{0, \dots, k-2\}$; $P_a$ denotes the orthogonal projection onto
$W_a = \mathrm{span}\{\chi_\xi : \vtwo(\xi) = a\}$, of dimension $d_a = 2^{k-2-a}$. The mean-zero
space is $V = W_0 \oplus \cdots \oplus W_{k-2}$, an orthogonal decomposition.
\end{definition}

%==============================================================================
\section{The spectral reduction}\label{sec:reduction}
%==============================================================================

\begin{proposition}\label{prop:reduction}
Let $Q = Q_k$ with $Q[a,b] = \|P_a U P_b\|_2$. Then
$|\lambda_2(T_k)| \le \rho(Q) \le \cert(k)$.
\end{proposition}

\begin{proof}
\emph{(i) Perron split.} Since $\mathbf{1}^{\!\top} T = \mathbf{1}^{\!\top}$, the mean-zero space
$V = \ker \mathbf{1}^{\!\top} = \mathbf{1}^\perp$ is $T$-invariant. In a basis adapted to
$\C^N = \mathrm{span}(\mathbf{1}) \oplus V$, $T$ is block-triangular
$\bigl(\begin{smallmatrix} 1 & 0 \\ d & A \end{smallmatrix}\bigr)$ with $A = T|_V$ (the vector
$d$, the $V$-component of $T\mathbf{1}$, is nonzero by Remark~\ref{rem:notds} but harmless). The
characteristic polynomial factors, so $\mathrm{spec}(T) = \{1\} \uplus \mathrm{spec}(A)$ and
$|\lambda_2(T)| \le \rho(A)$.

\emph{(ii) Adjoint and Gelfand.} $A^\ast = (P_V T P_V)^\ast|_V = P_V U P_V =: U_V$, the
compression of $U$ to $V$ ($V$ need not be $U$-invariant; only the compression is used), and
$\rho(A) = \rho(U_V)$. Gelfand's formula $\rho(U_V) \le \|U_V^p\|_2^{1/p}$ holds for every matrix
and every $p$; no normality or diagonalisability enters. This is the step that makes the argument
immune to the $10^k$ eigenvector condition numbers that invalidated the earlier perturbative
route.

\emph{(iii) Level majorisation.} For $g \in V$ let $\alpha(g)_a := \|P_a g\|_2$. Since
$P_a P_V = P_a$, the triangle inequality gives, for each level $a'$,
\[
  \|P_{a'} U_V g\| = \Bigl\| \sum_b P_{a'} U P_b g \Bigr\|
  \le \sum_b Q[a', b]\, \alpha(g)_b = (Q\, \alpha(g))_{a'},
\]
i.e.\ $\alpha(U_V g) \le Q\,\alpha(g)$ entrywise. As $Q \ge 0$ this iterates:
$\alpha(U_V^p g) \le Q^p \alpha(g)$. The levels are orthogonal and exhaust $V$, so
$\|U_V^p g\|_2 = \|\alpha(U_V^p g)\|_2 \le \|Q^p\|_2 \|g\|_2$, whence
$\rho(U_V) \le \lim_p \|Q^p\|^{1/p} = \rho(Q)$.

\emph{(iv) Weighted row sum.} For the diagonal similarity $S = \mathrm{diag}(2^0, \dots, 2^{k-2})$
and the nonnegative matrix $Q$,
$\rho(Q) = \rho(S Q S^{-1}) \le \|S Q S^{-1}\|_\infty = \max_a \sum_b Q[a,b]\, 2^{a-b} = \cert(k)$.
\end{proof}

%==============================================================================
\section{Block structure: three instances of one engine}\label{sec:blocks}
%==============================================================================

Throughout, $x$ is odd, $v(r) := \vtwo(3r+1) \ge 1$, and $3^{-1}$ denotes the inverse of 3 in the
relevant $\Z/2^m$.

\subsection{Coset uniformity and the splitting \texorpdfstring{$U = U_{\mathrm{clean}} + D$}{U = Uclean + D}}

\begin{lemma}[CU]\label{lem:CU}
Let $x$ be odd with $v := \vtwo(3x+1) < K$, and $q_0 := (3x+1)/2^v$ (odd). Then for every $m$,
\[
  \vtwo\bigl(3(x + m 2^K) + 1\bigr) = v, \qquad
  \Syr(x + m 2^K) = q_0 + 3m\,2^{K-v},
\]
and as $m$ ranges over $\Z/2^K$ the residue $\Syr(x + m 2^K) \bmod 2^K$ is uniform on the coset
$q_0 + 2^{K-v}(\Z/2^K)$ of size $2^v$, each value attained exactly $2^{K-v}$ times.
\end{lemma}

\begin{proof}
$3(x + m2^K) + 1 = 2^v q_0 + 3m 2^K = 2^v (q_0 + 3m\,2^{K-v})$, and the bracket is odd because
$q_0$ is odd and $K - v \ge 1$; this freezes the valuation and gives the affine formula. For the
count: $m \mapsto 3 \cdot 2^{K-v} m$ is an endomorphism of $\Z/2^K$ whose image is the subgroup
$2^{K-v}(\Z/2^K)$ of order $2^v$ (3 is a unit) and whose kernel has order $2^{K-v}$; every image
point has exactly $2^{K-v}$ preimages, and adding $q_0$ translates the image to the coset.
\end{proof}

Apply Lemma~\ref{lem:CU} with $K = k$. For every row $r$ with $v(r) < k$, averaging
$\chi_\eta$ over the uniform coset collapses to a geometric sum which is $2^v$ or $0$ according
to whether $2^v \mid \eta$, yielding the masked-phase formula
\begin{equation}\label{eq:masked}
  (U \chi_\eta)(r) = \mathbf{1}[\,v(r) \le \vtwo(\eta)\,]\; \omega^{\eta\, q(r)}, \qquad
  q(r) := (3r+1)/2^{v(r)}, \qquad (v(r) < k).
\end{equation}
Exactly one odd residue fails the hypothesis: $2^k \mid 3r+1$ has the unique odd solution
$r^\ast = -3^{-1} \bmod 2^k = (4^{\lceil k/2 \rceil} - 1)/3$, where $3r^\ast + 1 =
2^{2\lceil k/2 \rceil}$ exactly. Define $U_{\mathrm{clean}}$ by \eqref{eq:masked} with the
$r^\ast$-row set to zero, and $D := U - U_{\mathrm{clean}}$. Then:

\begin{itemize}
\item[\textbf{(R1)}] $P_a U_{\mathrm{clean}} P_b = 0$ for $a \ge b$ (strict upper-triangularity;
  a by-product of the survivor analysis below);
\item[\textbf{(R2)}] $D = e_{r^\ast} c^{\ast}$ is rank~1, where $c$ is the probability vector
  $c[t] = 2^{-k} \#\{m : \Syr(r^\ast + m2^k) \equiv t \bmod 2^k\}$; hence
  $\|P_a D P_b\| = u_a v_b$ with $u_a := \|P_a e_{r^\ast}\|$, $v_b := \|P_b c\|$;
\item[\textbf{(R3)}] $u_a = 2^{-(a+1)/2}$ exactly: any state unit vector has flat character
  energy $1/N$ per frequency, so $u_a^2 = d_a / N = 2^{-1-a}$.
\end{itemize}

By the triangle inequality, for all $a, b$:
\begin{equation}\label{eq:triangle}
  Q[a,b] \;\le\; \|P_a U_{\mathrm{clean}} P_b\| + u_a v_b.
\end{equation}

\subsection{The two remaining elementary inputs for Lemma A}

\begin{lemma}[SB: shell bijection]\label{lem:SB}
Fix $1 \le j \le k-2$. The map $\psi(r) := (3r+1)/2^j \bmod 2^{k-j}$ is a bijection from the
shell $S_j := \{ r \text{ odd}, 0 < r < 2^k,\ v(r) = j \}$ onto the odd residues modulo
$2^{k-j}$.
\end{lemma}

\begin{proof}
Injectivity: $\psi(r) = \psi(r')$ gives $3r + 1 \equiv 3r' + 1 \pmod{2^k}$, so $r = r'$ (3 a
unit, both in $[0, 2^k)$). Cardinality: $v(r) = j$ pins $r$ to a single (odd) residue class
modulo $2^{j+1}$, so $|S_j| = 2^{k-j-1}$, which equals the number of odd residues modulo
$2^{k-j}$; injective plus equal finite cardinality gives bijective, with multiplicity exactly
one.
\end{proof}

\begin{lemma}[S4: odd geometric sum]\label{lem:S4}
For $m \ge 1$ and $\Sodd(\alpha, m) := \sum_{q \text{ odd},\, 0 \le q < 2^m} \omega^{\alpha q}$:
if $2^{k-1} \mid \alpha$ then $|\Sodd| = 2^{m-1}$ (full); otherwise $\Sodd = 0$ if and only if
$k - m \le \vtwo(\alpha) \le k - 2$ (the dead band); below the band it is a nonzero partial sum.
\end{lemma}

\begin{proof}
Substitute $q = 2\ell + 1$ and sum the geometric series with ratio $\omega^{2\alpha}$.
\end{proof}

\subsection{Lemma A: the clean upper cascade is a scaled isometry}

\begin{theorem}[Lemma A]\label{thm:lemmaA}
For all $k$ and all $0 \le a < b \le k-2$, with $d := b - a$, the block
$B = P_a U_{\mathrm{clean}} P_b$ satisfies $B^\ast B = 2^{-d} I$; in particular
$\|P_a U_{\mathrm{clean}} P_b\|_2 = 2^{-d/2}$ exactly. Moreover $P_a U_{\mathrm{clean}} P_b = 0$
for $a \ge b$ (R1).
\end{theorem}

\begin{proof}[Proof (survivor analysis)]
The matrix entry of $U_{\mathrm{clean}}$ between source $\chi_\eta$ ($\vtwo(\eta) = b$) and
target $\chi_\xi$ ($\vtwo(\xi) = a$) is, up to the $1/N$ normalisation,
$\hat{U}(\eta, \xi) = \sum_{r,\, 1 \le v(r) \le b} \omega^{\eta q(r) - \xi r}$. Group by the
shell $j = v(r)$. On shell $j$, $r = 3^{-1}(2^j q - 1)$ with $q = q(r)$, so the phase is
$q\,\alpha_j + \xi 3^{-1}$ with $\alpha_j := \eta - \xi\, 2^j 3^{-1}$. The valuations are
\[
  \vtwo(\alpha_j) = a + j \ (j < d), \qquad \vtwo(\alpha_j) = b \ (j > d), \qquad
  \vtwo(\alpha_j) \ge b \ (j = d),
\]
and in every case $\vtwo(\alpha_j) \ge j$, so $\omega^{q \alpha_j}$ depends only on
$q \bmod 2^{k-j}$. By Lemma~\ref{lem:SB} the shell sum is exactly
$\omega^{\xi 3^{-1}} \Sodd(\alpha_j, k-j)$, and by Lemma~\ref{lem:S4} every shell lands in the
dead band $[j, k-2]$ \emph{except} possibly $j = d$, which survives if and only if
$\alpha_d \in \{0, 2^{k-1}\}$, with full modulus $2^{k-d-1}$. Counting: each $\eta$ has exactly
$2^d$ owners $\xi$ (two residue classes modulo $2^{k-d}$, each contributing $2^{d-1}$ values of
valuation $a$ in $[0, 2^{k-1})$), giving the diagonal
$(B^\ast B)[\eta, \eta] = N^{-2} \cdot 2^d \cdot (2^{k-d-1})^2 = 2^{-d}$; and for fixed $\xi$
exactly one of the two candidate sources lies in the index range $[0, 2^{k-1})$, so distinct
columns have disjoint row support and the off-diagonal entries vanish. For $a \ge b$ the same
computation gives $\vtwo(\alpha_j) = b \in [j, k-2]$ for every shell $j \le b$, so every shell
is annihilated: this is (R1).
\end{proof}

%==============================================================================
\section{Lemma B: the defect mass, by an integer collision count}\label{sec:lemmaB}
%==============================================================================

\begin{theorem}[Lemma B]\label{thm:lemmaB}
$\displaystyle \sum_{b=0}^{k-2} v_b^2 \;\le\; \|c\|_2^2 \;=\; \frac{\mathrm{coll}(k)}{4^k}
\;\le\; \frac{3 \cdot 2^k}{4^k}, \qquad\text{i.e.}\qquad \|v\|_2 \le \sqrt{3}\, 2^{-k/2},$
where $\mathrm{coll}(k) = \sum_t \mathrm{cf}[t]^2$ and
$\mathrm{cf}[t] = \#\{m \in [0,2^k) : \Syr(r^\ast + m 2^k) \equiv t \bmod 2^k\}$.
\end{theorem}

\begin{proof}[Proof sketch (full details in the repository)]
Since $3r^\ast + 1 = 2^{2\lceil k/2\rceil}$, the fiber is
$\Syr(r^\ast + m2^k) = \mathrm{oddpart}(a + 3m) \bmod 2^k$ with
$a = 2^{2\lceil k/2\rceil - k} \in \{1, 2\}$, an arithmetic progression of length $2^k$ in the
integers. Decompose it by 2-adic shells $A_j = \{x = a + 3m : \vtwo(x) = j\}$, of sizes
$|A_j| = 2^{k-1-j}$ ($j \le k-1$) plus one top atom.

\emph{FACT 1 (per-shell injectivity).} On each shell, $x \mapsto (x/2^j) \bmod 2^k$ is
injective: if $x = 2^j u$, $x' = 2^j u'$ with $u \equiv u' \pmod{2^k}$, then
$3(m - m') = 2^j(u - u') \equiv 0 \pmod{2^{j+k}}$; as $\gcd(2,3)=1$ this forces
$2^{j+k} \mid m - m'$, and $|m - m'| < 2^k$ gives $m = m'$. Hence the fiber count is $0/1$ per
shell, and $\mathrm{coll} = \sum_{j,j'} |R_j \cap R_{j'}|$ over the shell images $R_j$.

The diagonal contributes exactly $2^k$. The cross terms obey
$\sum_{j < j'} |R_j \cap R_{j'}| \le \sum_{j' \ge 2} (j'-1) 2^{k-1-j'} + (k-1) = 2^{k-1} - 1$
(the closed form includes the top atom's at most $k-1$ partners, covered by the truncation slack
of the geometric series). Hence $\mathrm{coll} \le 2^k + 2(2^{k-1} - 1) < 3 \cdot 2^k$. The
sharp constant is $\mathrm{coll}/2^k \to 31/12$, comfortably inside 3. Finally
$\sum_b v_b^2 = \|c\|^2 - 1/N \le \|c\|^2$, the dropped term being the Perron coefficient of the
probability vector $c$.
\end{proof}

%==============================================================================
\section{The assembly}\label{sec:assembly}
%==============================================================================

\begin{proof}[Proof of Theorem~\ref{thm:main}]
Fix a row $a \le k-2$ and set $e := k - a \ge 2$. Write the weight as
$2^{a-b} = 2^a (1/2)^b$. By \eqref{eq:triangle}, Theorem~\ref{thm:lemmaA} and (R1),
\[
  \sum_b Q[a,b]\, 2^{a-b}
  \;\le\; \underbrace{\sum_{d=1}^{e-2} 2^{-d/2}\, 2^{-d}}_{\text{clean cascade, truncated}}
  \;+\; \underbrace{u_a 2^a \sum_{b=0}^{k-2} v_b\, 2^{-b}}_{\text{defect, all } b}.
\]
Two remarks. The clean sum \emph{truncates} at $d \le e-2$ because $b \le k-2$; at the top row
$a = k - 2$ it is empty. The defect is summed over \emph{all} $b$, which absorbs its
contribution to the upper blocks with no separate estimate.

For the defect factor, Cauchy--Schwarz against the geometric weight and Theorem~\ref{thm:lemmaB}
give, with no information about the profile of $(v_b)$,
\[
  \sum_{b=0}^{k-2} v_b\, 2^{-b}
  \;\le\; \|v\|_2 \Bigl( \sum_{b \ge 0} 4^{-b} \Bigr)^{1/2}
  \;\le\; \sqrt{3}\, 2^{-k/2} \cdot \tfrac{2}{\sqrt{3}} \;=\; 2^{1 - k/2},
\]
and with $u_a 2^a = 2^{(a-1)/2}$ (R3) the defect term is at most
$2^{(a-1)/2}\, 2^{1-k/2} = 2^{-(e-1)/2}$. Hence
\[
  \sum_b Q[a,b]\, 2^{a-b} \;\le\; f(e) \;:=\; \sum_{d=1}^{e-2} 2^{-3d/2} + 2^{-(e-1)/2}.
\]
The difference $f(e+1) - f(e) = 2^{-(e-1)/2}\bigl(2^{-(e-1)} - (1 - 2^{-1/2})\bigr)$ is positive
only at $e = 2$, so $f$ increases from $f(2) = 2^{-1/2}$ to its maximum
$f(3) = 2^{-3/2} + 2^{-1} = 0.8535\ldots$ and decreases thereafter toward
$G_{\mathrm{up}} = (2^{3/2} - 1)^{-1} = 0.5469\ldots$. Therefore
$\cert(k) \le f(3) < 1$ for every $k \ge 3$, and Proposition~\ref{prop:reduction} completes the
proof.
\end{proof}

\begin{remark}[Why truncation matters]\label{rem:truncation}
Without the truncation one would bound the row sum by
$G_{\mathrm{up}} + 2^{-1/2} = 1.25 > 1$ and wrongly conclude that the $L^2$ defect bound is
insufficient (an earlier version of the assembly did exactly this and introduced the per-level
Lemma C as a necessity). Row-wise, the two terms of $f$ trade off: rows with a long cascade have
a small defect weight and vice versa, with the worst case at $e = 3$.
\end{remark}

\begin{remark}[Sharpness and measured values]\label{rem:sharp}
The envelope is asymptotically sharp at the bottom row ($e = k$), where both the true row sum and
$f(k)$ converge to $G_{\mathrm{up}}$; the row that sets the constant, $e = 3$, has a wide margin
(true value $\approx 0.63$ against $f(3) = 0.854$). Ground-truth values (dense operator,
both parities; reproducible with \texttt{fable\_assembly\_check.py}):

\begin{center}
\begin{tabular}{lcccccc}
\toprule
$k$ & 4 & 6 & 8 & 10 & 12 & 13 \\
\midrule
$\cert(k)$ & 0.6117 & 0.6338 & 0.6347 & 0.6345 & 0.6345 & 0.6344 \\
$\rho(Q_k)$ & 0.5020 & 0.5535 & 0.5661 & 0.5673 & 0.5656 & 0.5644 \\
$\|c\|_2\, 2^{k/2}$ & 1.6202 & 1.6105 & 1.6081 & 1.6075 & 1.6073 & 1.6072 \\
\bottomrule
\end{tabular}
\end{center}

\noindent All row bounds $R_a \le f(k-a)$ hold with margin at every tested $k$, and the rank-1
factorisation $Q[a,b] = u_a v_b$ on $a \ge b$ is machine-exact.
\end{remark}

%==============================================================================
\section{The per-level sharpening (Lemma C)}\label{sec:lemmaC}
%==============================================================================

Theorem~\ref{thm:main} needs only the $L^2$ mass of the defect. The per-level structure is
nevertheless elementary and sharp, and improves the constant.

\begin{theorem}[Lemma C]\label{thm:lemmaC}
$v_b \le \tfrac{3}{4}\, 2^{-b}\, 2^{-k/2}$ for all $k$ and $0 \le b \le k-2$.
\end{theorem}

The proof (in the repository) rests on a frequency-doubling identity for the fiber exponential
sum $S_k(\xi) = \sum_t \mathrm{cf}[t]\, e^{-2\pi i \xi t / 2^k}$:

\begin{lemma}[H: homometry]\label{lem:H}
For $k \ge 2$ and $2\xi \notin \{0, 2^{k-1}\} \bmod 2^k$: $S_k(2\xi) = S_{k-1}(\xi)$ exactly.
\end{lemma}

\begin{proof}[Proof idea]
Split the fiber sum by the parity of $m$. One parity makes $3m + a$ odd and contributes a
geometric series that vanishes away from the two excluded frequencies; the other parity halves
the scale and reproduces the $(k-1)$-fiber exactly (the parity flip swaps the offsets
$a = 1 \leftrightarrow 2$, matching the parity of $k-1$).
\end{proof}

Iterating Lemma~\ref{lem:H} reduces the level-$b$ energy to a single scale $p = k - b$, where a
Parseval identity converts it into the half-shift autocorrelation $\mathrm{coll}(p) - A_p$, which
the shell decomposition of Theorem~\ref{thm:lemmaB} bounds by $3 \cdot 2^{p-1} - 2$ (the $j = 0$
shell cancels identically from the autocorrelation). Feeding Theorem~\ref{thm:lemmaC} through
the truncated envelope of Section~\ref{sec:assembly} replaces $2^{-(e-1)/2}$ by $2^{-(e+1)/2}$
and yields
\[
  \cert(k) \;\le\; \max_{e \ge 2}\Bigl[ \sum_{d=1}^{e-2} 2^{-3d/2} + 2^{-(e+1)/2} \Bigr]
  \;=\; 0.6559\ldots \qquad (\text{maximum at } e = 4),
\]
remarkably close to the measured $0.6345$.

%==============================================================================
\section{Lean formalisation}\label{sec:lean}
%==============================================================================

Theorem~\ref{thm:main} is machine-checked in Lean~4 against Mathlib (v4.27.0), end to end and
with no undischarged hypothesis. The endpoint is \texttt{gap\_certificate\_unconditional}: for
every $k \ge 3$, every eigenvalue $\mu \ne 1$ of the operator $T_k$ --- \emph{defined inside the
formalisation} from the lift window, not axiomatised --- satisfies $\|\mu\| < 0.853554$.

The project is 18 files and 449 theorem/lemma declarations, with no \texttt{sorry}, no
\texttt{native\_decide}, and every \texttt{\#print axioms} set a subset of
$\{\texttt{propext}, \texttt{Classical.choice}, \texttt{Quot.sound}\}$. (That check is a subset
test, not a string match: a number of declarations legitimately use fewer axioms.) The
correspondence with this paper is:

\begin{itemize}
\item \textbf{$T_k$ itself} --- \texttt{TransferOperator.lean}. \texttt{Tcount k u r} counts the
  lifts $m < 2^k$ of the odd residue $r$ with $\mathrm{Syr}(r + m 2^k) \equiv u$, which is
  Section~\ref{sec:setup}'s definition transcribed. Nothing downstream is about an abstract
  operator satisfying axioms.
\item \textbf{Proposition~\ref{prop:reduction}} (the spectral reduction) ---
  \texttt{OperatorChain.lean}: the Perron split on $\ker \mathbf{1}^\top$, the compression, and
  the Gelfand step. The one-sided form is load-bearing and is the form proved: $T_k$ is
  \emph{not} doubly stochastic and its stationary vector is \emph{not} uniform.
\item \textbf{Lemma~\ref{lem:CU}} (coset uniformity), including the fiber count, and the
  counting halves of Lemma~\ref{lem:SB} --- \texttt{CountingLemmas.lean}.
\item \textbf{Theorem~\ref{thm:lemmaB}} (the collision count) --- \texttt{CollisionBound.lean},
  reaching the assembly through \texttt{Assembly.hL2\_discharged}.
\item \textbf{Theorem~\ref{thm:lemmaA}} --- the survivor analysis of
  Section~\ref{sec:blocks} in \texttt{LemmaA.lean} (the $S4$ dead band, the valuation table, the
  survivor rule), the entry theorem and $B^\ast B = 2^{-d} I$ in \texttt{GramIdentity.lean}
  (\texttt{gram\_upper}), the owner count and the ownership partition in
  \texttt{OwnerCount.lean} and \texttt{OwnerPartition.lean}, and (R1), the $a \ge b$ vanishing,
  as an operator identity in \texttt{OperatorBlock.lean}.
\item \textbf{The assembly of Section~\ref{sec:assembly}} --- \texttt{Assembly.lean}, with the
  envelope $f(e) \le f(3) < 1$ and the row-sum bound; instantiated at the concrete $T_k$ in
  \texttt{ManifestInstance.lean}.
\item \textbf{The sign-blindness of Section~\ref{sec:scope}} --- \texttt{SignBlind.lean}. This
  file is a \emph{sibling}, not a link in the chain above: it imports nothing beyond
  \texttt{CountingLemmas.lean}, and \texttt{gap\_certificate\_unconditional} neither uses nor
  needs it. It proves that for every odd $c$ and every valuation pattern with
  $\sum j_i + 1 \le k$ the count is $2^{\,k-1-\sum j_i}$, hence the probability $2^{-\sum j_i}$
  with no $k$ in it, hence identical finite-dimensional distributions for $3x+1$ and $3x-1$.
  What it does \emph{not} do is construct a measure on $\Z_2$; see the caveat below.
\end{itemize}

\begin{remark}[What the sign-blindness formalisation does not contain]
\texttt{SignBlind.lean} establishes the combinatorial content of the equality of the two
valuation-process laws --- every cylinder mass, at every depth, inside the band
$\sum j_i + 1 \le k$. It does not establish that equality as a statement about measures on
$\Z_2$. Two steps are missing, and they are not equally hard.

The passage from cylinder masses to a measure on $\N^\N$ is the lighter one, and requires no new
theory: Mathlib already carries the product measure \texttt{Measure.infinitePiNat} indexed by
$\N$, its projective-limit property, the identification of its value on measurable cylinders
with the finite product, and the independence of its coordinates. The target law is an exact fit
for \texttt{geometricPMF} at parameter $1/2$ shifted by one, since
$(1/2)^n (1/2) = 2^{-(n+1)}$. What \texttt{SignBlind.lean} computes \emph{is} the cylinder mass,
so the assembly is bounded work against existing library material, and no general Kolmogorov
extension theorem enters.

The harder step is tying the counting measure to Haar measure on $\Z_2$ --- that Haar restricted
to level $k$ is uniform on residues mod $2^k$, the balls being cosets of $2^k \Z_2$. This is
standard mathematics, but Mathlib provides \texttt{PadicInt} without the accompanying Haar
development, so it is the piece that would need new material. It is avoidable at the cost of
stating the result at finite level, where no $\Z_2$ appears at all.

We are explicit about this because the distance between ``all cylinder masses agree'' and ``the
laws are equal'' is exactly the kind of step it is tempting to elide.
\end{remark}

Three formalisation dividends are worth recording. FACT~1, parametrised by the lift index, needs
no modulo-3 argument at all. The identity $(2s^2)^{a+1} = 1$ for $s = 2^{-1/2}$ keeps every
exponent in $\N$, so the formal assembly involves no real exponentiation. And the dead band of
Lemma~\ref{lem:S4} must be stated in divisibility form rather than as
$k - m \le \vtwo(\alpha) \le k-2$: the valuation form is false at $\alpha = 0$ under the
convention $\vtwo(0) = 0$, and the formalisation caught this where inspection had not.

The formalisation followed this paper's own proof of Theorem~\ref{thm:lemmaA} without
alteration, including the step that had been expected to be hardest. The off-diagonal entries of
$B^\ast B$ do not require a cancellation among roots of unity; as the proof above says, distinct
columns have disjoint row support, so every off-diagonal summand already carries a zero factor.

\textbf{Not formalised.} Lemma C (Section~\ref{sec:lemmaC}) --- it sharpens the constant rather
than feeding it, so Theorem~\ref{thm:main} does not depend on it --- and the non-degeneracy
$g_c \ne 0$, which the certificate also does not use. Nothing else in the chain of
Theorem~\ref{thm:main} is paper-only.

%==============================================================================
\section{What this theorem does and does not say}\label{sec:scope}
%==============================================================================

\subsection*{What the theorem is a statement about}

This subsection is placed first because the natural reading of Theorem~\ref{thm:main} --- ``the
mod-$2^k$ Collatz chain mixes uniformly fast'' --- is wrong, and the correct reading is more
interesting.

Let $K_k$ be the \emph{honest} 2-adic kernel: the transition matrix on odd residues mod $2^k$
obtained by conditioning Haar measure on the \emph{full} 2-adic fibre of a residue, rather than
on the finite lift window $[0,2^k)$. Then $K_k$ is not merely fast-mixing, it is
\textbf{nilpotent off the Perron eigenvalue}: $|\lambda_2(K_k)| = 0$, so the chain reaches
stationarity in finitely many steps. This follows from the closed form of the Bernstein--Lagarias
conjugacy \cite{BernsteinLagarias1996} --- on the shift side each Syracuse step deletes at least
one low bit, so the non-Perron part is a nilpotent shift --- and is not proved here. (Numerically
$|\lambda_2(K_k)|$ computed with $2^{18}$ lifts is $\approx 10^{-2}$ and falls steadily as the
fibre is resolved, consistent with $0$.)

So \emph{the mixing of the averaged mod-$2^k$ Collatz dynamics is trivial}, and
$\cert(k)$ is not measuring it. What $T_k$ differs from $K_k$ by is exactly the substitution of
the finite lift window for the full fibre, and that difference is supported on one row:

\[
  \operatorname{rank}(T_k - K_k) = 1 \quad\text{for every } k,
  \qquad \|T_k - K_k\|_1 \approx 0.25\text{--}0.50 .
\]

Both verified for $k = 3, \dots, 8$. \textbf{The defect is rank one but it is not small}, and it
does not shrink as $k$ grows. A rank-one perturbation of that size can in principle drag an
eigenvalue to the unit circle and destroy mixing altogether.

\textbf{Theorem~\ref{thm:main} says that it does not, at any scale.} That is the honest content:
replacing the full 2-adic fibre average by the finite lift window --- which is what every
computation on the mod-$2^k$ Collatz chain actually does --- perturbs a nilpotent kernel into one
whose sub-dominant spectrum still sits below $0.853553\ldots$, uniformly in $k$, with the true
values near $0.27$. It is a \emph{uniform spectral stability statement for the finite-window
approximation}, and its audience is anyone computing mod-$2^k$ Collatz statistics over a finite
window and relying, usually silently, on that approximation not degenerating as the window
refines.

\textbf{One limit, stated plainly: this is a statement about the spectrum, not about mixing
times.} $T_k$ is markedly non-normal and its eigenvector condition number grows like $10^k$
(Section~\ref{sec:intro-history}, item (i)), so $|\lambda_2| \le 0.8536$ does \emph{not} convert
into a total-variation mixing bound uniform in $k$. Spectrum and mixing come apart for this
family. Conflating them is precisely the error of the withdrawn March 2026 draft, and it is not
repaired by the present proof --- only avoided, because the present proof never needs
eigenvectors.

It says nothing about individual deterministic orbits.

\textbf{The withdrawn cycle claim.} An earlier version of this project asserted that the uniform
gap eliminates non-trivial Collatz cycles. The inference is false, by a control experiment: the
$3x - 1$ map (equivalently, Collatz on negative integers) has non-trivial cycles --- in Syracuse
form exactly $\{5, 7\}$ and $\{17, 25, 37, 55, 41, 61, 91\}$, besides the trivial $\{1\}$, and no
others below $2 \times 10^6$ --- yet
its transfer operator passes the identical certificate, with an even stronger measured gap. The
averaging over $2^k$ lifts that defines $T_k$ erases exactly the deterministic orbit structure a
cycle lives in, in the same way that the doubling map is mixing while having dense periodic
orbits. Ruling out Collatz cycles requires instruments sensitive to specific orbits (heights,
linear forms in logarithms), not the spectrum of the averaged operator.

We claim no novelty for the principle behind this control, only for the specific certificate it
defeats. That the sign is invisible to this class of instrument is recorded in the literature:
Kontorovich and Lagarias observe that the differences between the dynamics of $3x+1$ and $5x+1$
on the integers are ``invisible at the level of measure theory'' \cite{KontorovichLagarias2010},
and note that the conjugacy results of Bernstein and Lagarias \cite{BernsteinLagarias1996} carry
the identical ergodic theory to all $ax+b$ maps --- which includes $3x-1$. Matthews states the
same phenomenon from the other direction, that for relatively prime type maps the prediction of
overall cycling or divergence is independent of the remainder values \cite{Matthews2010}. That
$3x-1$ has non-trivial cycles is folklore. What is ours here is narrower and sufficient: the
observation that \emph{this} certificate, with \emph{these} constants, is passed by $3x-1$ with a
stronger measured gap, which is what makes the withdrawn inference demonstrably false rather than
merely unsupported.

Likewise the theorem has
no bearing on divergent trajectories; the determinism-versus-probability gap quantified in the
project's audit (a large-deviation shortfall of roughly a factor 12 for every counting-based
route) remains the actual wall, consistent with the state of the art after \cite{Tao2019}.

What the theorem does provide: a natural family of arithmetically structured, non-normal,
non-doubly-stochastic operators of exponentially growing dimension with a uniform, explicit,
elementarily provable spectral gap, whose sharp structure ($B^\ast B = 2^{-d} I$ blocks, a rank-1
defect at the base-4 repunit $r^\ast$, the $31/12$ collision constant) is completely understood.

\subsection*{What this is good for}

Three uses, none of which is progress on Collatz.

\textbf{A falsification test for averaged arguments.} The control experiment above generalises
into a cheap and decisive check that costs a referee ten minutes. \emph{Run the candidate
argument against $3x - 1$.} If it goes through unchanged, it is wrong, because $3x-1$ has the
cycles listed above. This works because averaged models are typically blind to the sign of the
shift, and that blindness is not a heuristic: for every odd $c$ and every valuation pattern
$(j_1, \dots, j_t)$ with $\sum j_i + 1 \leq k$, the number of odd residues mod $2^k$ realising
that pattern under $x \mapsto 3x + c$ is $2^{\,k-1-\sum j_i}$, independent of $c$. Equivalently
the induced probability is $2^{-\sum j_i}$, with no $k$ in it, so the two valuation processes
have identical finite-dimensional distributions. This is machine-checked (\texttt{SignBlind.lean},
sorry-free); the underlying conjugacy is classical \cite{BernsteinLagarias1996, Lagarias1985}.

\textbf{Reusable formalisation.} The Lean development is public domain and modular. Independently
of Collatz it supplies 2-adic fibre counts for $x \mapsto 3x + c$ mod $2^k$, an orthonormal
character basis on $\Z/2^k$ with projections graded by valuation, and a Gram identity
$B^\ast B = 2^{-d} I$ established without forming an adjoint.

\textbf{A negative result.} Knowing which families of argument are closed is worth time in a
problem that has consumed a great deal of it. Every route surveyed here that makes contact with
the conjecture exits through the same door --- an orbit-quantified statement about a hypothetical
infinite trajectory --- and no averaged spectral datum reaches that door.

We claim no application outside mathematics and are not aware of one.

\subsection*{Reproducibility}

All claims are reproducible from
\url{https://github.com/m4cd4r4/collatz-conjecture-spectral-gap}: the proof documents
(\texttt{THEOREM.md} and the lemma files), the numerical ground-truth checks
(\texttt{fable\_assembly\_check.py}, \texttt{verify\_assembly.py},
\texttt{audit\_halfshift\_s4.py}, \texttt{lemmaB\_fact1\_rigorous.py}), the $3x{-}1$ control
(\texttt{probe\_cycle\_link.py}, \texttt{CYCLE\_CLAIM\_REFUTED.md}), and the Lean project
(\texttt{lean/}).

\begin{thebibliography}{9}

\bibitem{Lagarias2010}
J.~C. Lagarias (ed.), \emph{The Ultimate Challenge: The $3x+1$ Problem}, American Mathematical
Society, 2010.

\bibitem{Tao2019}
T.~Tao, \emph{Almost all orbits of the Collatz map attain almost bounded values}, Forum of
Mathematics, Pi 10 (2022), e12. arXiv:1909.03562.

\bibitem{BernsteinLagarias1996}
D.~J. Bernstein and J.~C. Lagarias, \emph{The $3x+1$ conjugacy map}, Canadian Journal of
Mathematics 48 (1996), 1154--1169.

\bibitem{Lagarias1985}
J.~C. Lagarias, \emph{The $3x+1$ problem and its generalizations}, American Mathematical Monthly
92 (1985), 3--23.

\bibitem{KontorovichLagarias2010}
A.~V. Kontorovich and J.~C. Lagarias, \emph{Stochastic models for the $3x+1$ and $5x+1$
problems}, in \cite{Lagarias2010}. arXiv:0910.1944.

\bibitem{Matthews2010}
K.~R. Matthews, \emph{Generalized $3x+1$ mappings: Markov chains and ergodic theory}, in
\cite{Lagarias2010}.

\bibitem{Mori2024}
L.~Mori, \emph{$C^\ast(T_1, T_2)$ irreducible on $\ell^2(\N)$ iff Collatz}, arXiv:2411.08084
(2024).

\bibitem{mathlib}
The mathlib Community, \emph{The Lean mathematical library}, CPP 2020.

\end{thebibliography}

\end{document}
